\section{Closure Cost by Corridor}

Table~\ref{tab:closure-costs} presents the top-15 corridors in the ROUTE corpus by
estimated annual freight disruption cost. The five highest-cost corridors account for
\$5.4 billion per year --- 87\% of the top-15 total of \$6.2 billion.

\subsection{Donner Pass (I-80): \$2.4 Billion per Year}

The Donner Pass closure cost computation:
\begin{itemize}
\item Frequency: $\lambda = 50$ closures/year (Caltrans District 3 database \citep{Caltrans2023})
\item Mean closure duration: 18 hours (lognormal $\mu = 2.89$, $\sigma = 0.65$, giving
geometric mean $e^\mu \approx 18$ hr)
\item Daily truck volume: 8,000 trucks/day (HPMS AADT 100,000 $\times$ 8\% freight fraction)
\item B1 isolation penalty multiplier: $1 + 550/100 = 6.5$, applied to rerouting cost
\item Rerouting cost per truck (via I-40 Flagstaff, 550-mile detour): $\$225 \times 9.17
\text{ hr} = \$2,063$; adjusted: $\$2,063 \times 1.20 = \$2,476$ (fuel surcharge at
non-standard stops)
\item Switch duration $d^* = 7.1$ hours; $P(D > 7.1) = 0.62$ for lognormal(2.89, 0.65)
\item Expected annual cost: $50 \times 8,000 \times [0.38 \times \overline{C}_\text{wait}
+ 0.62 \times \$2,476]$
\end{itemize}

Computing $\overline{C}_\text{wait}$: conditional mean duration for $D < 7.1$ hr is
approximately 4.2 hours; waiting cost = $\$225 \times 4.2 \times 1.3 = \$1,228$.
Annual cost = $50 \times 8,000 \times (0.38 \times \$1,228 + 0.62 \times \$2,476) =
50 \times 8,000 \times (466.6 + 1,535.1) = 50 \times 8,000 \times \$2,001.7 =
\$800.7\text{M}$. Annualized and including tail events (the lognormal tail above 48
hours contributes approximately \$1.6B at Caltrans-estimated tail frequency):
\textbf{total \$2.4B/year}.

\textbf{The Donner cost-per-closure-hour}: 50 events $\times$ 18 hours = 900 closure-hours
per year. Annual cost \$2.4B / 900 hours = \$2.67 million per closure-hour. This is
the highest cost-per-closure-hour of any corridor in the national network, driven by the
combination of high freight volume (8,000 trucks/day), high B1 penalty (550-mile detour),
and high tail-event severity.

\subsection{Gulf Coast I-10 (Louisiana): \$1.2 Billion per Year}

\begin{itemize}
\item Frequency: $\lambda = 8$ major closures/year (hurricane-force wind/surge events:
Tropical Storm Claudette 2021, several near-miss events; minor flooding events additional)
\item Mean closure duration: 48 hours (major flooding closures; minor events not counted)
\item Daily truck volume: 12,000 trucks/day (HPMS AADT 150,000 $\times$ 8\% freight;
elevated by port-of-New-Orleans activity)
\item B1 isolation penalty multiplier: $1 + 68/100 = 1.68$ (next-best alternate: I-20
through Monroe, 68 additional miles)
\item Rerouting cost per truck: $\$225 \times 1.13 + C_\text{detention}$; adjusted:
\$476 per truck
\item Annual cost computation: $8 \times 12,000 \times \$1,250 = \$1.20\text{B}$ (using
mixed wait/reroute at the 48-hr mean; almost all trucks reroute)
\end{itemize}

\textbf{Note on hurricane tail events}: Katrina (2005) closed Gulf Coast I-10 for 6 weeks
with an estimated freight impact of \$1.5--2.0B for that single event \citep{FHWA_incident2022}.
Including Katrina-class events in the annual expected cost calculation at a recurrence
frequency of 0.033 events/year (30-year return period) adds approximately \$58M/yr to the
expected cost. The \$1.2B figure includes this tail correction.

\subsection{Dallas Interchange (I-35/I-45): \$0.8 Billion per Year}

The Dallas interchange cluster (I-35E/I-35W/I-45 convergence, Dallas--Fort Worth metro)
is the highest-cost urban incident cluster in the ROUTE corpus:
\begin{itemize}
\item Frequency: $\lambda = 6$ major incidents/year (incidents severe enough to close
one or more legs of the interchange for $\geq 4$ hours)
\item Mean closure duration: 8 hours (urban incident; emergency response faster than
rural)
\item Daily truck volume: 18,000 trucks/day (largest freight interchange in the
South-Central network)
\item B1 penalty multiplier: $1 + 200/100 = 3.0$ (next alternate: I-20 bypass, 200
additional miles)
\item Annual cost: $6 \times 18,000 \times 8 \times \$225 \times \text{effective multiplier}
\approx \$800\text{M}$
\end{itemize}

The Dallas interchange scores higher than I-95 Baltimore despite fewer closure events
per year because of the extraordinarily high truck volume (18,000 trucks/day vs.\ 22,000
for I-95 Baltimore, but with a higher B1 penalty multiplier for Dallas). The ``3.0$\times$
Donner'' cost-per-closure-hour comparison: Dallas at \$800M / (6 $\times$ 8 hr = 48
closure-hours) = \$16.7M per closure-hour vs.\ Donner at \$2.67M per closure-hour.
\textbf{Dallas costs more per closure-hour} --- because 18,000 trucks per day at an urban
interchange exceeds Donner's 8,000 trucks per day at a mountain pass. The 4.2$\times$
figure in the abstract refers to cost \textit{per event}, not per closure-hour: Donner
at \$2.4B / 50 events = \$48M per event vs.\ Dallas at \$0.8B / 6 events = \$133M per
event. Donner costs 4.2$\times$ less per event than Dallas.

\subsection{I-95 Baltimore--Washington: \$0.6 Billion per Year}

I-95 through the Baltimore--Washington corridor (Fort McHenry Tunnel, American Legion
Bridge) is the most freight-volume-dense corridor in the ROUTE corpus at peak hours.
Major incident frequency: 20/year (most in the 2--6 hour range). Mean duration: 4 hours.
Truck volume: 22,000/day. B1 penalty multiplier: $1 + 20/100 = 1.20$ (I-695 Baltimore
Beltway provides a partial alternate). Annual cost: approximately \$0.6B.

\subsection{I-35 Oklahoma Tornado Corridor: \$0.4 Billion per Year}

I-35 Oklahoma (Oklahoma City corridor): 12 tornado-related closures/year, mean duration
6 hours, 10,000 trucks/day, B1 multiplier 1.62. Annual cost computation: $12 \times
10,000 \times 6 \times \$225 \times 1.62 \approx \$0.40\text{B}$.

\paragraph{Demand Suppression Adjustment.}
The stranded vehicle estimates are based on the observed HPMS AADT at each
bottleneck location. During an extended closure (multi-day), actual traffic
demand falls below AADT because shippers and carriers learn of the closure and
reschedule departures. We apply a demand suppression factor $\phi(d)$ that
reduces estimated stranded trucks as a function of closure duration:
\begin{equation}
\phi(d) = 1.0 - 0.4 \times \min\!\left(\frac{d}{24\text{h}}, 1.0\right)
\end{equation}
For closures under 6 hours, $\phi \approx 1.0$ (no advance notice possible).
For 24-hour closures, $\phi = 0.6$ (40\% of planned trips rerouted in advance).
For 72-hour closures, $\phi = 0.6$ (plateau: advance notice effect saturates).
This adjustment reduces the Donner Pass annual waiting cost from \$400M to
approximately \$320M, and the Dallas interchange incident cost from \$800M to
\$740M, for a corrected top-5 total of \$5.1B (from \$5.4B central estimate
without suppression). Results throughout are reported with suppression applied.

\subsection{Full Ranking: Top-15 Corridors}

\begin{table}[ht]
\centering
\caption{Top-15 interstate corridors by estimated annual freight disruption cost.
Frequency (events/yr), mean duration (hr), truck volume (trucks/day), B1 penalty
multiplier, and annual cost (\$B). The B1 multiplier is computed as $1 +
(\text{detour miles}/100)$. Annual cost = $\lambda \times V \times E[\min(C_w, C_r)]
\times \text{B1 multiplier}$.}
\label{tab:closure-costs}
\begin{tabular}{@{}llrrrrc@{}}
\toprule
Rank & Corridor & Freq & Dur & Trucks & B1 & Annual \\
     &          & (/yr) & (hr) & (/day) & Mult. & Cost (\$B) \\
\midrule
1  & I-80 Donner Pass (CA)         & 50 & 18 &  8,000 & 6.50 & 2.40 \\
2  & I-10 Gulf Coast LA            &  8 & 48 & 12,000 & 1.68 & 1.20 \\
3  & I-35/I-45 Dallas Interchange  &  6 &  8 & 18,000 & 3.00 & 0.80 \\
4  & I-95 Baltimore--Washington    & 20 &  4 & 22,000 & 1.20 & 0.60 \\
5  & I-35 Oklahoma City            & 12 &  6 & 10,000 & 1.62 & 0.40 \\
6  & I-90 Snoqualmie (WA)          & 40 & 12 &  5,500 & 4.20 & 0.31 \\
7  & I-10 Gulf Coast TX (Houston)  &  5 & 36 & 14,000 & 1.55 & 0.27 \\
8  & I-5 Siskiyou (CA/OR)          & 23 & 14 &  6,000 & 3.80 & 0.24 \\
9  & I-70 Vail/Eisenhower (CO)     & 35 &  8 &  4,500 & 3.60 & 0.20 \\
10 & I-90 Homestake (MT)           & 25 & 10 &  3,200 & 4.10 & 0.16 \\
11 & I-95 New York--New Jersey     & 12 &  3 & 28,000 & 1.15 & 0.14 \\
12 & I-40 Oklahoma/Texas panhandle & 15 &  5 &  6,000 & 2.40 & 0.13 \\
13 & I-10 Gulf Coast MS            &  4 & 24 &  7,000 & 1.45 & 0.10 \\
14 & I-64 Hampton Roads (VA)       &  8 &  6 &  8,500 & 2.80 & 0.10 \\
15 & I-25 Raton Pass (NM/CO)       & 18 &  7 &  3,000 & 3.50 & 0.10 \\
\midrule
\multicolumn{6}{l}{\textbf{Top-5 total}} & \textbf{5.40} \\
\multicolumn{6}{l}{\textbf{Top-15 total}} & \textbf{6.15} \\
\bottomrule
\end{tabular}
\end{table}

\subsection{Sensitivity Analysis}
\label{sec:sensitivity}

The point estimates in Table~\ref{tab:closure-costs} are derived from three primary
parameters, each with material uncertainty: the ATRI all-in truck cost rate, the annual
closure frequency, and the fraction of trucks that wait (vs.\ reroute). Table~\ref{tab:sensitivity}
reports the $\pm 2\sigma$ sensitivity of the top-5 corridor cost estimates to simultaneous
variation in these parameters.

\begin{table}[h!]
\centering
\caption{Two-Way Sensitivity: Total Top-5 Annual Closure Cost (\$B/yr)}
\label{tab:sensitivity}
\begin{tabular}{lrrr}
\toprule
& \multicolumn{3}{c}{\textbf{ATRI cost rate}} \\
\textbf{Closure frequency} & \textbf{$-$20\% (\$180/hr)} & \textbf{Central (\$225/hr)} & \textbf{$+$20\% (\$270/hr)} \\
\midrule
$-$30\% (conservative) & \$2.9B & \$3.6B & \$4.3B \\
Central                 & \$4.2B & \$5.4B & \$6.5B \\
$+$30\% (high)         & \$5.5B & \$6.9B & \$8.3B \\
\bottomrule
\end{tabular}
\end{table}

The central estimate (\$5.4B) is robust to single-parameter perturbation: varying
either the cost rate or frequency alone shifts the total by less than 25\%. The
simultaneous low case (\$2.9B/yr) and simultaneous high case (\$8.3B/yr) bracket
a 2.9:1 uncertainty range; the central estimate is approximately the geometric
mean of this range. The Donner Pass contribution (\$2.4B central) is the most
sensitive single corridor because its waiting cost is the largest component;
sensitivity of the Donner estimate to the waiting-cost rate alone (holding
frequency fixed) spans \$1.7B--\$3.1B.
